\documentclass[12pt]{article}
\usepackage[margin=1in]{geometry}
\usepackage{amssymb,amsmath}
\usepackage{booktabs}
\usepackage{graphicx}
\usepackage{caption}
\usepackage[round]{natbib}
\usepackage{hyperref}
\hypersetup{colorlinks=true, linkcolor=black, citecolor=black, urlcolor=blue}
\usepackage{setspace}
\onehalfspacing

\title{\textbf{Generative-AI Exposure in Sub-Saharan African 
Labour Markets}}
\author{Abiodun Adetokunbo\thanks{Augustine University, Lagos, Nigeria. Email:abiodun.adetokunbo@augustineuniversity.edu.ng}}
\date{\today}

\begin{document}
\maketitle

\begin{abstract}
\noindent This study provides a measurement of occupational exposure
 to generative AI across workforce of four Sub-Saharan African economies, 
 using harmonised microdata for 23,444 workers from LSMS-ISA surveys in Nigeria, 
 Tanzania, Uganda, and Ethiopia. Exposure is 
 concentrated in a thin layer of the labour force such that the median worker faces 
 near-zero exposure. Two thirds of workers are in agricultural or elementary
 occupations whose exposure is below 0.06, and the most exposed tenth of workers
 accounts for roughly a third of all exposure. The dividing line is sectoral 
 rather than contractual as agricultural workers, who make up 59 percent of the 
 pooled workforce, have a mean exposure of 0.05, while non-agricultural wage and 
 self-employed workers are equally exposed at 0.14, formal and informal wage 
 workers do not differ. Across countries, in a 
 sample of 32 African economies, country-level AI exposure falls steeply with 
 the informal employment rate and the agricultural self-employment share, and 
 rises with income per capita. This study finds that exposure is not merely lower on 
 average but structurally confined to a small, urban, white-collar minority.

\vspace{6pt}
\noindent\textbf{Keywords:} generative AI, occupational exposure, informality, structural transformation, Sub-Saharan Africa, LSMS-ISA.

\noindent\textbf{JEL codes:} J21, J24, O33, O55.
\end{abstract}

\section{Introduction}

The rapid diffusion of generative AI has prompted a wave of estimates of how exposed 
jobs are to the new technology \citep{felten2021,eloundou2024,gmyrek2023}. 
Most of this evidence concerns high-income labour markets, and the cross-country 
work that does cover developing economies builds occupational exposure from 
adult-skills surveys such as PIAAC and STEP \citep{pizzinelli2023,cazzaniga2024}. 
Those instruments, by design, sample the formal and largely urban workforce. 
They do not represent the subsistence farmers, own-account traders, and casual 
labourers who make up the majority of employment in low-income Sub-Saharan Africa. 
As a result, the question of who in these economies is actually exposed, and 
how exposure is distributed across occupational structure, has not been 
answered with representative data.

This paper takes a different data source as its starting point. 
The Living Standards Measurement Study Integrated Surveys on Agriculture (LSMS-ISA) 
are nationally representative household surveys that record the main-job occupation 
of every working member, including the agricultural and informal 
majority \citep{lsmsisa}. The paper assembles the LSMS-ISA wave with coded 
occupations for four countries, Nigeria (2023), Tanzania (2020), Uganda (2019), 
and Ethiopia (2021), and map each worker's occupation to a task-level measure of 
generative-AI exposure through a transparent occupation crosswalk. 
The pooled sample contains 23,444 workers.

This study is descriptive and its contribution is in three parts. First, it measures 
the level 
and distribution of AI exposure across a representative Sub-Saharan workforce and finds 
exposure is low and highly concentrated such that the pooled mean is 0.088 
on a zero-to-one automation scale, the median worker is somewhat unexposed, and 
the most exposed ten percent of workers carry about a third of total exposure. 
Second, it shows that the concentration is occupational and sectoral rather than 
contractual. The gap is between agriculture and everything else, not between formal 
and informal jobs such that agricultural workers average 0.05, while non-agricultural wage 
employees and own-account workers are both at 0.14, and formality among wage workers 
makes no difference. Third, it shows that the within-country pattern scales up. 
In a cross-section of 32 African economies, aggregate exposure falls with the 
informal employment rate and with agricultural self-employment, and rises with 
income, such that the four countries observed at the worker level sit on a regional gradient.

These results qualify the view that developing countries are
less exposed to AI \citep{pizzinelli2023,cazzaniga2024}. Using data that include 
the informal and agricultural workforce, it finds that exposure in these economies 
is not simply lower on average but is structurally confined to a small, urban, 
white-collar minority, with the large agricultural and own-account majority 
insulated by the nature of their work. This is a measurement and mapping
contribution as the study do not estimate the effect of AI on
employment.

\subsection{Related literature}

The paper connects to three strands of work. The first measures the exposure of
jobs to automation and to AI. The task framework of \citet{autor2003} and
\citet{acemoglu2019,acemoglu2022} treats an occupation as a bundle of tasks and
asks which of those tasks a technology can perform. A series of occupation-level
indices apply this idea, from the early computerisation probabilities of
\citet{frey2017} to the AI occupational exposure scores of \citet{felten2021} and
\citet{webb2020} and the large-language-model exposure of \citet{eloundou2024}.
The exposure measure used here follows this tradition, aggregating task-level
capability ratings for generative models to occupations \citep{althoff2026}.
A parallel literature estimates the realised effects of generative AI on
productivity and on labour demand, mostly in high-income or online settings
\citep{brynjolfsson2023,noy2023,hui2024}. This paper does not estimate effects but instead maps exposure.

The second strand carries exposure measurement across countries and raises the
concern that AI may widen gaps between nations \citep{alonso2020}. Cross-country
estimates consistently find that lower-income economies are less exposed, because
they employ fewer workers in the professional and clerical occupations that AI
reaches \citep{pizzinelli2023,cazzaniga2024,gmyrek2023}. These estimates rest on
adult-skills surveys such as PIAAC and STEP that sample the formal workforce in a
limited set of countries, so they describe the exposed minority rather than the
whole labour force. This paper uses household survey data that include the
agricultural and informal majority, and asks how exposure is distributed once that
majority is in the sample.

The third strand concerns informality and structural transformation. Informal and
own-account employment dominates the labour market in low-income economies
\citep{laporta2014}, and a large share of work remains in low-productivity
agriculture \citep{gollin2014}. These features are usually studied as obstacles to
development. Here they appear instead as the reason AI exposure is so concentrated,
since the agricultural and own-account work that absorbs most of the workforce is
made up of tasks current models cannot perform.

\section{Data and measurement}

\subsection{Surveys and occupation coding}

The study uses LSMS-ISA wave that records coded occupations for each of 
four countries, the Nigeria General Household Survey-Panel Wave 5 (2023), 
the Tanzania National Panel Survey Round 5 (2020), the Uganda National 
Panel Survey (2019), and the Ethiopia Socioeconomic Panel Survey Wave 5 (2021). 
Each survey records the main-job occupation of working household members. Nigeria, 
Tanzania, and Uganda code occupation at the four-digit ISCO-08 level \citep{isco08}; 
Ethiopia records occupation at the one-digit major-group level. Occupation coverage 
is complete among workers in Nigeria and Uganda. In Tanzania the labour 
module codes occupation only for non-farm work, so the study adds subsistence farmers, 
identified from the agricultural-labour question, at the agricultural occupation 
level. The Ethiopian sample is smaller and selected, and treats its level with 
caution and rely on it mainly for the major-group profile.

\subsection{Exposure measure and crosswalk}

The study measures occupational exposure to generative AI with task-level capability 
ratings that score each O*NET task for whether current generative models can 
perform it \citep{althoff2026}. Ratings are aggregated to occupations using task 
weights, yielding an automation-exposure score on a zero-to-one scale for 
each O*NET-SOC occupation. In order to attach this score to survey workers, 
the study builds a crosswalk from each survey's occupation classification to O*NET-SOC: 
ISCO-88 and ISCO-08 to the 2010 Standard Occupational Classification, and then to 
O*NET-SOC, using the standard correspondence tables. Where one source occupation 
maps to several O*NET occupations I take the unweighted mean. 
I also compute exposure from an independent occupation-level automation index and this implies 
the same cross-occupational ranking, and conclusions concern that 
ranking rather than the cardinal scale.

\subsection{Harmonised variables}

For each worker, I record sex, age, urban or rural residence, 
the four-digit occupation and its ISCO major group, and an employment-status 
category. Status is anchored on the occupation: agricultural workers are 
identified from ISCO major group 6 together with agricultural elementary 
occupations, and the remaining workers are split into wage employees and 
own-account or self-employed workers using each survey's main-job employer question. 
For wage workers, I additionally construct an informality flag from the presence of 
a written contract, pension contributions, or income-tax withholding, depending on 
what each survey records. The resulting pooled file covers 23,444 workers
(7,708 in Nigeria, 7,481 in Tanzania, 7,671 in Uganda, and 584 in Ethiopia).
Table~\ref{tab:summary} describes the sample. The four workforces differ in the
ways the structural account predicts: the agricultural share ranges from 28
percent in the more urban Nigerian sample to 78 percent in Tanzania, and mean
exposure moves inversely with it.

\begin{table}[ht]
\centering
\caption{Sample composition by country}
\label{tab:summary}
\small
\begin{tabular}{lrrrrrr}
\toprule
Country & $N$ & \% female & Mean age & \% urban & \% agriculture & Mean exposure \\
\midrule
Nigeria 2023  & 7{,}708 & 50.5 & 26.0 & 19.0 & 28.4 & 0.122 \\
Tanzania 2020 & 7{,}481 & 50.5 & 39.3 & 3.5  & 77.9 & 0.083 \\
Uganda 2019   & 7{,}671 & 65.2 & 26.6 & 22.4 & 72.8 & 0.055 \\
Ethiopia 2021 & 584     & 38.0 & 26.5 & 44.5 & 22.4 & 0.144 \\
\midrule
All (pooled)  & 23{,}444 & 55.0 & 30.5 & 15.8 & 58.6 & 0.088 \\
\bottomrule
\end{tabular}

\vspace{4pt}
\begin{minipage}{0.9\linewidth}
\footnotesize
Workers with a coded main-job occupation. Mean exposure is the task-weighted
automation-exposure score on a zero-to-one scale. The agriculture column is the
share of workers in agricultural occupations.
\end{minipage}
\end{table}

\section{The structure of AI exposure}

Generative-AI exposure reaches only a small part 
of the workforce in the sample. The pooled mean automation-exposure score is 0.088 
and the median is 0.076. One third of workers (33.7 percent) have near-zero 
exposure below 0.05, and two thirds (66.1 percent) are in agricultural or 
elementary occupations. Exposure is also concentrated among the few who have it as 
the most exposed ten percent of workers account for 36 percent of total exposure 
in the pooled sample.

Table~\ref{tab:major} reports employment shares and mean exposure 
for each ISCO major group. The exposure gradient is steep and follows a somewhat 
white-collar ordering. Clerical, technical, and professional occupations have the 
highest exposure, with clerical workers at 0.605, but together they are a small 
share of employment. The two largest groups, skilled 
agriculture (54 percent of workers) and elementary occupations (12 percent), 
have the lowest exposure, at 0.054 and 0.030. The combination of a 
steep occupational gradient and an employment distribution that is heavily 
weighted toward its low-exposure end is what produces a low and concentrated 
aggregate.

\begin{table}[ht]
\centering
\caption{Employment share and mean AI exposure by ISCO major group 
(pooled, 23,444 workers)}
\label{tab:major}
\begin{tabular}{lrr}
\toprule
ISCO major group & Share of workers & Mean exposure \\
\midrule
Clerical support             & 0.7\%  & 0.605 \\
Technicians and associates   & 1.8\%  & 0.361 \\
Managers                     & 5.1\%  & 0.210 \\
Professionals                & 4.6\%  & 0.172 \\
Service and sales            & 12.6\% & 0.143 \\
Plant and machine operators  & 3.6\%  & 0.143 \\
Craft and trades             & 5.4\%  & 0.051 \\
Skilled agriculture          & 54.0\% & 0.054 \\
Elementary occupations       & 12.1\% & 0.030 \\
\bottomrule
\end{tabular}

\vspace{4pt}
\begin{minipage}{0.85\linewidth}
\footnotesize
Exposure is the task-weighted share of an occupation's tasks that current generative models 
can perform, on a zero-to-one scale. Shares are of all workers with a coded occupation in 
the pooled four-country sample.
\end{minipage}
\end{table}

Table~\ref{tab:status} reports 
exposure by employment status. Agricultural workers, who are 58.6 percent of the pooled 
workforce, have a mean exposure of 0.050. Non-agricultural workers are more exposed, but 
within the non-agricultural sector, wage employees and own-account workers are exposed to 
exactly the same degree, both at 0.142. The insulating margin is therefore agriculture versus 
the rest of the economy.

\begin{table}[ht]
\centering
\caption{Employment share and mean AI exposure by employment status (pooled)}
\label{tab:status}
\begin{tabular}{lrr}
\toprule
Employment status & Share of workers & Mean exposure \\
\midrule
Agriculture                       & 58.6\% & 0.050 \\
Self-employed, non-agriculture    & 36.0\% & 0.142 \\
Wage employee, non-agriculture    & 3.5\%  & 0.142 \\
Other non-agriculture             & 1.9\%  & 0.163 \\
\bottomrule
\end{tabular}

\vspace{4pt}
\begin{minipage}{0.85\linewidth}
\footnotesize
Status is anchored on the occupation: 
agriculture is read from the ISCO major group, and remaining workers are 
split by the survey's main-job employer question. The ``other'' category 
is the Ethiopian non-agricultural workers, for whom only the one-digit occupation 
is available.
\end{minipage}
\end{table}

Among non-agricultural wage workers, 
for whom formality can be measured, informal workers (no written contract, pension, 
or tax withholding) have a mean exposure of 0.112 and formal workers 0.100. 
If anything informal wage workers are marginally more exposed, because informal 
wage work in these economies includes clerical and sales activity that generative 
models can perform. Geography and sex add little as urban workers average 0.095 
against 0.087 for rural workers, and men 0.092 against 0.085 for women. 
The concentration of exposure is occupational, and the occupational structure, 
not residence or contract type, is what places most workers out of reach of the 
technology.

\section{From within-country to cross-country}
 
Table~\ref{tab:country} reports, for each of the four micro countries, 
the mean worker-level exposure measured and an independent country-level 
measure of AI substitution exposure drawn from a 32-country cross-section of 
African economies. The worker-level levels are not perfectly comparable across 
countries, because the Ethiopian occupation code is coarse and its sample selected, 
which inflates its measured level. The country-level measure, which does not share 
that selection, places Ethiopia among the least exposed, consistent with its very 
high agricultural share.

\begin{table}[ht]
\centering
\caption{Worker-level and country-level AI exposure, four micro-data countries}
\label{tab:country}
\begin{tabular}{lrr}
\toprule
Country & Worker-level mean & Country-level exposure \\
\midrule
Nigeria 2023   & 0.122 & 0.333 \\
Ethiopia 2021  & 0.144 & 0.221 \\
Tanzania 2020  & 0.083 & 0.238 \\
Uganda 2019    & 0.055 & 0.214 \\
\bottomrule
\end{tabular}

\vspace{4pt}
\begin{minipage}{0.85\linewidth}
\footnotesize
The worker-level mean is the pooled automation-exposure score from the LSMS sample. 
The country-level measure is an independent index of occupational substitution 
exposure aggregated to the national level.
\end{minipage}
\end{table}

Across the full 32-country cross-section, the same insulation logic that 
operates within countries operates between them. Country-level AI exposure 
falls with the informal employment rate (Spearman rank correlation of $-0.74$, 
28 countries) and with the agricultural self-employment 
share ($-0.75$, 30 countries), and rises with income per 
capita ($+0.83$, 31 countries). Figure~\ref{fig:macro} plots exposure against 
the informal employment rate and marks the four countries observed at the worker 
level as they lie along the regional gradient. The cross-country relationship is 
descriptive, but it confirms that the structural feature shown with microdata, 
the confinement of exposure to a small non-agricultural minority, is what drives 
the well-known result that poorer and more informal economies appear less exposed.

\begin{figure}[ht]
\centering
\includegraphics[width=0.8\linewidth]{figures/fig_macro_exposure_informality}
\caption{Country-level AI exposure against the informal employment rate, 
32 African economies. The four LSMS micro-data countries are marked.}
\label{fig:macro}
\end{figure}

\section{Limitations}

First, the exposure measure is built on 
US task content and is applied to African occupations through a crosswalk, therefore
it captures the task composition of an occupation, not local differences in how 
that occupation is performed, so it is best read as a measure of an occupation's 
intrinsic exposure rather than of realised local risk. Second, cross-country 
comparisons of the cardinal level are imperfect. Ethiopian occupation is coded 
only at the major-group level, and the Tanzanian occupation module originally 
omitted subsistence farmers, which was added at the agricultural level. 
The conclusions rest on the occupational gradient and on the pooled distribution, 
rather than on the precise national means. 
Third, the analysis is descriptive. The study measures exposure, not impact, and a
companion comparison of a Nigerian worker panel between 2018 and 2023 shows no
exposure-related change in employment over a window in which generative AI had
barely been adopted, so this cross-section is best read as a pre-impact baseline
rather than as evidence of an effect.

\section{Conclusion}

Using LSMS-ISA microdata that include the agricultural 
and informal workforce, the study provides a measurement of generative-AI exposure 
across the occupational structure of four Sub-Saharan African economies.
Exposure is low, steeply graded across occupations, and concentrated in a small 
white-collar minority. The agricultural and own-account majority is insulated by 
the nature of its work rather than by the informality of its contracts. 
The same structure scales to a 32-country regional gradient in which more 
informal and more agricultural economies are less exposed. For policy and for 
measurement, the implication is that aggregate exposure statistics understate how 
unequally the technology reaches workers in these economies as the gains and risks of 
generative AI, for now, are confined to the thin formal-sector top of the labour
market, while the structural transformation that would expose the majority is lagging.

\bibliographystyle{plainnat}
\bibliography{refs}

\end{document}
